% Source coverage: physical PDF pp. 200--205 (printed pp. 177--182), beginning at the Section 19.3 heading on p. 200 and ending with the section's final paragraph above the Section 19.4 heading on p. 205.
% Numbered equations: (19.3.1)--(19.3.20), with all unnumbered displays. Figures/tables: none. Footnotes: 1. Inline chapter references: [7], [8]. Uncertainties: none.

\subsection{Spontaneously Broken Approximate Symmetries}

In the previous section we dealt with exact symmetries of the action that do not leave the vacuum invariant, and are said to be spontaneously broken. We shall now consider the effect of adding small symmetry-breaking terms to the action in such a theory. Such spontaneously broken approximate symmetries are important in the theory of strong interactions, and in some areas of condensed matter physics. As we shall see, the spontaneous breakdown of an approximate symmetry does not lead to the appearance of massless Goldstone bosons, but of low-mass spinless particles, often called \emph{pseudo-Goldstone bosons}~\hyperref[ch19-ref-7]{[7]}.

We continue here to treat translationally invariant theories, in which the effective action is expressed as in Eq.~\eqref{eq:19.2.8} in terms of an effective potential $V(\phi)$ that depends on a set of spacetime-independent scalar field expectation values $\phi_n$. For an action that obeys some set of approximate continuous symmetries with generators $t_\alpha$, the effective potential may be written
\begin{equation}
V(\phi)=V_0(\phi)+V_1(\phi),
\tag{19.3.1}\label{eq:19.3.1}
\end{equation}
where $V_0(\phi)$ satisfies the invariance condition\footnote{We denote general symmetry generators as $t_\alpha$, $t_\beta$, etc., in contrast with the independent broken symmetry generators, for which the subscript $\alpha$ takes values $a$, $b$, etc.}
\begin{equation}
\sum_{mn}\frac{\partial V_0(\phi)}{\partial\phi_n}(t_\alpha)_{nm}\phi_m=0
\tag{19.3.2}\label{eq:19.3.2}
\end{equation}
and $V_1(\phi)$ is a small correction due to the symmetry breaking in the action. Suppose that this perturbation shifts the minimum of the potential from $\phi_0$, the minimum of $V_0(\phi)$, to $\bar\phi=\phi_0+\phi_1$, where $\phi_1$ is small, of first order in the symmetry-breaking perturbation. The equilibrium condition for the vacuum is then
\begin{equation}
\left.\frac{\partial V(\phi)}{\partial\phi_n}\right|_{\phi=\phi_0+\phi_1}=0.
\tag{19.3.3}\label{eq:19.3.3}
\end{equation}
The zeroth-order term on the left is just
$[\partial V_0(\phi)/\partial\phi_n]_{\phi=\phi_0}$, which vanishes because $\phi_0$ is defined as the minimum of $V_0(\phi)$. Thus the first-order terms must also vanish, and so:
\begin{equation}
\left.\sum_m\frac{\partial^2 V_0(\phi)}
{\partial\phi_n\partial\phi_m}\right|_{\phi=\phi_0}\phi_{1m}
+\left.\frac{\partial V_1(\phi)}{\partial\phi_n}\right|_{\phi=\phi_0}=0.
\tag{19.3.4}\label{eq:19.3.4}
\end{equation}
Eq.~\eqref{eq:19.2.6} holds here with $V$ replaced with the invariant term $V_0$ and $\bar\phi$ replaced with $\phi_0$:
\begin{equation}
\left.\sum_n\frac{\partial^2 V_0(\phi)}
{\partial\phi_n\partial\phi_m}\right|_{\phi=\phi_0}
(t_\alpha)_{n\ell}\phi_{0\ell}=0.
\tag{19.3.5}\label{eq:19.3.5}
\end{equation}
Hence multiplying Eq.~\eqref{eq:19.3.4} with $(t_\alpha\phi_0)_n$ and summing over $n$ gives
\begin{equation}
\left.(t_\alpha\phi_0)_n
\frac{\partial V_1(\phi)}{\partial\phi_n}\right|_{\phi=\phi_0}=0.
\tag{19.3.6}\label{eq:19.3.6}
\end{equation}
Recalling the interpretation of $V(\phi)$ as the generating function for one-particle-irreducible graphs, and noting that in the absence of the perturbation $V_1$ the Goldstone components of $\phi_n$ are those in the direction of $t_\alpha\phi_0$ (see Eq.~\eqref{eq:19.2.33}), the left-hand side of Eq.~\eqref{eq:19.3.6} is proportional to the sum of all `tadpole' graphs, in which a pseudo-Goldstone boson disappears into the vacuum. Eq.~\eqref{eq:19.3.6} may thus be paraphrased as the condition that, to first order in $V_1$, \emph{pseudo-Goldstone bosons have no tadpoles}.

The moral of this calculation is that if we do not start with a zeroth-order vacuum expectation value that satisfies Eq.~\eqref{eq:19.3.6}, then even a small perturbation will produce a large change in $\bar\phi$, invalidating the expansion of $\bar\phi$ around $\phi_0$. Fortunately, for compact Lie groups it is always possible to choose $\phi_0$ to satisfy Eq.~\eqref{eq:19.3.6}. To see this, note that the invariance of the potential $V_0(\phi)$ under a group of real linear transformations $\phi\to L\phi$ implies that if $\phi_*$ is one minimum of the potential, then so is $L\phi_*$. For continuous groups of transformations, we can always parameterize the transformations as $L(\theta)$, in such a way that
\begin{equation}
\frac{\partial L(\theta)}{\partial\theta^\alpha}L^{-1}(\theta)
=iN_{\alpha\beta}(\theta)t_\beta,
\tag{19.3.7}\label{eq:19.3.7}
\end{equation}
where $N_{\alpha\beta}$ is a non-singular matrix depending on the group parameters $\theta_\alpha$. (Recall that in a real representation, it is $it_\alpha$ rather than $t_\alpha$ that is real.) Now consider the function $V_1(L(\theta)\phi_*)$. Where the group is compact, $L(\theta)\phi_*$ maps out a compact manifold as $\theta$ runs over the group volume, and as long as $V_1(\phi)$ is continuous, it must have a minimum on any such compact surface, say at a point $L(\theta_*)\phi_*$. The derivative of $V_1(L(\theta)\phi_*)$ with respect to $\theta_\alpha$ is
\begin{equation}
\frac{\partial V_1(L(\theta)\phi_*)}{\partial\theta_\alpha}
=\left.\sum_n\frac{\partial V_1(\phi)}{\partial\phi_n}
\right|_{\phi=L(\theta)\phi_*}
N_{\alpha\beta}(\theta)
\bigl(it_\beta L(\theta)\phi_*\bigr)_n.
\tag{19.3.8}\label{eq:19.3.8}
\end{equation}
This must vanish at the minimum $\theta_*$, and since $N_{\alpha\beta}$ is non-singular, this implies that
\begin{equation}
0=\left.\sum_n\frac{\partial V_1(\phi)}{\partial\phi_n}
\right|_{\phi=L(\theta_*)\phi_*}
\bigl(t_\beta L(\theta_*)\phi_*\bigr)_n.
\tag{19.3.9}\label{eq:19.3.9}
\end{equation}
But then Eq.~\eqref{eq:19.3.6} is satisfied if we make the choice $\phi_0=L(\theta_*)\phi_*$.

Eq.~\eqref{eq:19.3.6} is known as a \emph{vacuum alignment} condition~\hyperref[ch19-ref-8]{[8]}, because it generally has the effect of forcing the direction of the symmetry breaking by the vacuum into some sort of alignment with the symmetry-breaking terms in the Hamiltonian. For instance, consider the case of $SO(N)$ spontaneously broken to $SO(N-1)$, introduced in the previous section. In the absence of any symmetry-breaking perturbation, there is no way to tell \emph{which} $SO(N-1)$ subgroup is left unbroken; if the dynamics of the theory leads to a ground state that is invariant under the $SO(N-1)$ subgroup of $SO(N)$ that leaves some $N$-vector $\phi_{0n}$ invariant, then by performing an $SO(N)$ rotation we can find a ground state that is invariant under the $SO(N-1)$ subgroup that leaves any other $N$-vector invariant. If we add a perturbation that transforms under $SO(N)$ like, say, the component $\sum_n u_n\phi_n$ of an $N$-vector $\phi_n$ (not necessarily consisting of elementary scalars), then the Hamiltonian is invariant under a specific $SO(N-1)$ subgroup of $SO(N)$, consisting of rotations that leave the vector $u$ invariant. Without a vacuum alignment condition, we might think that the remaining exact symmetry is $SO(N-2)$, consisting of those rotations that leave invariant \emph{both} $u$ and the vector $\phi_0$ characterizing the vacuum symmetry. But with $V_1(\phi)=\sum_n u_n\phi_n$, the condition \eqref{eq:19.3.6} tells us that at the true vacuum, $\sum_n(t_\alpha\phi_0)_n u_n=0$ for all $SO(N)$ generators $t_\alpha$. The $SO(N)$ generators $t_\alpha$ span the space of all antisymmetric $N\times N$ matrices, so this condition requires that $\phi_0$ must be in the same direction as $u$, and so the unbroken symmetry is $SO(N-1)$, not $SO(N-2)$.

According to the general results of Section 16.1, the mass matrix $M_{ab}$ of the pseudo-Goldstone bosons is given to first order by
\begin{equation}
M_{ab}^2=\left.\sum_{mn}Z_{an}Z_{bm}
\frac{\partial^2 V(\phi)}{\partial\phi_m\partial\phi_n}
\right|_{\phi=\phi_0+\phi_1},
\tag{19.3.10}\label{eq:19.3.10}
\end{equation}
where $Z_{an}$ is the field renormalization constant defined by Eq.~\eqref{eq:19.2.39}. Since the mass matrix \eqref{eq:19.3.10} vanishes in zeroth order, the first-order terms give
\begin{equation}
M_{ab}^2=\sum_{mn}Z_{an}Z_{bm}
\left[
\left.\frac{\partial^3 V_0(\phi)}
{\partial\phi_\ell\partial\phi_m\partial\phi_n}
\right|_{\phi=\phi_0}\phi_{1\ell}
+\left.\frac{\partial^2 V_1(\phi)}
{\partial\phi_m\partial\phi_n}\right|_{\phi=\phi_0}
\right],
\tag{19.3.11}\label{eq:19.3.11}
\end{equation}
where $Z_{an}$ is here given by the zeroth-order approximation to \eqref{eq:19.2.40}:
\begin{equation}
Z_{an}=\sum_b F_{ab}^{-1}(it_b\phi_0)_n.
\tag{19.3.12}\label{eq:19.3.12}
\end{equation}
To calculate the mass matrix \eqref{eq:19.3.11}, we take $t$ in Eq.~\eqref{eq:19.2.55} to be one of the broken symmetry generators $t_a$, set $\phi=\phi_0$, and contract with $(t_b\phi_0)_m\phi_{1\ell}$:
\begin{align*}
0={}&
\left.\sum_{nm\ell}
\frac{\partial^3 V_0(\phi)}
{\partial\phi_\ell\partial\phi_m\partial\phi_n}
\right|_{\phi=\phi_0}
\phi_{1\ell}(t_a\phi_0)_n(t_b\phi_0)_m
\\
&+\left.\sum_{nm}
\frac{\partial^2 V_0(\phi)}
{\partial\phi_n\partial\phi_m}
\right|_{\phi=\phi_0}
(t_a\phi_1)_n(t_b\phi_0)_m
\\
&+\left.\sum_{n\ell}
\frac{\partial^2 V_0(\phi)}
{\partial\phi_n\partial\phi_\ell}
\right|_{\phi=\phi_0}
(t_at_b\phi_0)_n\phi_{1\ell}.
\end{align*}
The second term on the right-hand side vanishes according to Eq.~\eqref{eq:19.2.6}, while the third may be rewritten using Eq.~\eqref{eq:19.3.4}, leaving us with
\begin{equation}
\left.\sum_{nm\ell}
\frac{\partial^3 V_0(\phi)}
{\partial\phi_\ell\partial\phi_m\partial\phi_n}
\right|_{\phi=\phi_0}
\phi_{1\ell}(t_a\phi_0)_n(t_b\phi_0)_m
=
\left.\frac{\partial V_1(\phi)}{\partial\phi_n}
\right|_{\phi=\phi_0}
(t_at_b\phi_0)_n.
\tag{19.3.13}\label{eq:19.3.13}
\end{equation}
Using this in Eq.~\eqref{eq:19.3.11} then yields a formula for the pseudo-Goldstone boson mass matrix in terms of $V_1$:
\begin{align}
M_{cd}^2=-\sum_{ab}F_{ca}^{-1}F_{db}^{-1}
\Biggl[
&(t_a\phi_0)_n(t_b\phi_0)_m
\left.\frac{\partial^2 V_1(\phi)}
{\partial\phi_m\partial\phi_n}\right|_{\phi=\phi_0}
\notag\\
&+(t_at_b\phi_0)_n
\left.\frac{\partial V_1(\phi)}{\partial\phi_n}
\right|_{\phi=\phi_0}
\Biggr].
\tag{19.3.14}\label{eq:19.3.14}
\end{align}

For this to be a sensible mass matrix, it had better be positive. To see that it is, it is convenient to rewrite this result in terms of derivatives with respect to the group parameters $\theta_\alpha$. Differentiating Eq.~\eqref{eq:19.3.8} with respect to $\theta_\beta$, setting $\theta=\theta_*$, and using Eq.~\eqref{eq:19.3.9} and $\phi_0=L(\theta_*)\phi_*$ gives
\begin{equation}
M_{ab}^2=\sum_{cd\alpha\beta}
N_{a\alpha}^{-1}(\theta_*)N_{b\beta}^{-1}(\theta_*)
F_{ac}^{-1}F_{bd}^{-1}
\left.
\frac{\partial^2 V_1(L(\theta)\phi_*)}
{\partial\theta_\alpha\partial\theta_\beta}
\right|_{\theta=\theta_*}.
\tag{19.3.15}\label{eq:19.3.15}
\end{equation}
The matrix on the right is positive, because $\theta_*$ is the \emph{minimum} of the function $V_1(L(\theta)\phi_*)$.

This formula has a somewhat more familiar version, in terms of the vacuum expectation value of a double commutator of symmetry generators with the symmetry-breaking perturbation. Suppose that the symmetry-breaking perturbation $H_1$ in the Hamiltonian is a linear combination
\begin{equation}
H_1=\sum_n u_n\Phi_n
\tag{19.3.16}\label{eq:19.3.16}
\end{equation}
of operators $\Phi_n$ (not necessarily elementary scalar fields), which furnish a representation of the symmetry group with generators $t_\alpha$, in the sense that
\begin{equation}
[T_\alpha,\Phi_n]=-(t_\alpha)_{nm}\Phi_m,
\tag{19.3.17}\label{eq:19.3.17}
\end{equation}
where $T_\alpha$ are the quantum mechanical generators of the symmetry group. According to the results of Section 16.3, the symmetry-breaking part of the potential is
\begin{equation}
V_1(\phi)=\langle H_1\rangle_{\langle\Phi\rangle=\phi}
=\sum_n u_n\phi_n,
\tag{19.3.18}\label{eq:19.3.18}
\end{equation}
the subscript in the middle expression indicating that the expectation value is to be taken in the state of minimum energy in which $\Phi_n$ has expectation value $\phi_n$. The vacuum alignment condition \eqref{eq:19.3.6} then reads
\[
0=\sum_n u_n(t_\alpha\phi_0)_n,
\]
or using Eq.~\eqref{eq:19.3.17}
\begin{equation}
0=\langle[T_\alpha,H_1]\rangle_0,
\tag{19.3.19}\label{eq:19.3.19}
\end{equation}
the subscript $0$ now indicating that the expectation value is to be taken in the vacuum state, in which $\Phi_n$ has expectation value $\phi_{0n}$. Also, Eq.~\eqref{eq:19.3.14} gives the mass matrix here as
\[
M_{cd}^2=-\sum_{ab}F_{ca}^{-1}F_{db}^{-1}
\sum_n u_n(t_at_b\phi_0)_n,
\]
and, using Eq.~\eqref{eq:19.3.17}, this is
\begin{equation}
M_{cd}^2=-\sum_{ab}F_{ca}^{-1}F_{db}^{-1}
\langle[T_a,[T_b,H_1]]\rangle_0.
\tag{19.3.20}\label{eq:19.3.20}
\end{equation}
This is symmetric in $c$ and $d$. To see this, note that the Jacobi identity and group commutation relations may be used to write the difference of Eq.~\eqref{eq:19.3.20} and the same with $c$ and $d$ interchanged as a linear combination of terms $\langle[T_\alpha,H_1]\rangle_0$, which vanish according to the vacuum alignment condition \eqref{eq:19.3.19}. The mass matrix \eqref{eq:19.3.20} is also positive, because the point $\theta=0$ is at the \emph{minimum} of the vacuum energy $\langle\exp(-i\theta_aT_a)H_1\exp(i\theta_aT_a)\rangle_0$ for rotated vacuum states $\exp(i\theta_aT_a)\ket{\Omega}$.
